\begin{figure}[tb]
\centering
\begin{tikzpicture}
\begin{axis}[diagnostic axis,title={Even d flux},ylabel={$\psi_d/\psi_*$},
 xmin=-1.1,xmax=1.1,ymin=.99,ymax=1.04]
 \addplot[strong line,reference quantity,domain=-1:1,samples=81]{1+.02*x^2};
 \addplot[only marks,estimated quantity] coordinates {(-.8,1.0128) (.8,1.0128)};
 \draw[alternative pattern,guide quantity] (axis cs:-.8,1.0128)--(axis cs:.8,1.0128);
\end{axis}
\begin{axis}[diagnostic axis,at={(.5\textwidth,0)},anchor=south west,
 title={Odd q flux},ylabel={$\psi_q/\psi_*$},xmin=-1.1,xmax=1.1,ymin=-1.1,ymax=1.1]
 \addplot[strong line,derived quantity,domain=-1:1,samples=81]{x-.08*x^3};
 \addplot[only marks,estimated quantity] coordinates {(-.8,-.75904) (.8,.75904)};
 \draw[alternative pattern,guide quantity] (axis cs:-.8,-.75904)--(axis cs:.8,.75904);
\end{axis}
\end{tikzpicture}
\caption{A nonlinear symmetric example at $i_d=0$. Reflected q currents
give equal d flux and opposite q flux despite curved responses. Curves are
from the co-energy later specified in \cref{eq:synthetic-energy}; estimated quantity marks
identify a reflected pair.}
\label{fig:magnetic-symmetry}
\end{figure}
